% GENERATED FILE. Edit canonical lessons, then run npm run generate:books.
% canonical-source-hash: fnv1a64:8054d358a61deb47
% canonical-lessons: ZH-C475-shouxin, ZH-C475-xigai, ZH-C475-houbei, ZH-C475-shengao, ZH-C475-wan-bend, ZH-R475-review-470-475

\chapter{Head to Knees}
\label{ch:zh-head-to-knees}
\hlchaptermodality{475}

\begin{chapteropening}
\textbf{By the end of this chapter:} \emph{I can talk about the palm of my hand, my knees and back, how tall I am, and bending an arm.}

Complete the last lesson of chapter 475: I can talk about the palm of my hand, my knees and back, how tall I am, and bending an arm.
\end{chapteropening}

\section[\texorpdfstring{\zh{手心} (shǒuxīn)}{手心 (shǒuxīn)}]{\zh{手心} (shǒuxīn) — the palm of the hand}

\label{lesson:ZH-C475-shouxin}



\begin{quote}
Before the new one: say the Mandarin for to put one's strength into it, then the Mandarin for a finger.
\end{quote}

\subsection*{You'll want to know: \zh{手心}}
\textbf{\zh{手心}} — \emph{shǒuxīn} — \textquotedblleft{}the palm of the hand\textquotedblright{}.

\begin{quote}
\textbf{\zh{我手心很热是不是体温高了}} \emph{wǒ shǒuxīn hěn rè shì bu shì tǐwēn gāo le} — My palms are hot. Is my temperature up?
\end{quote}

\begin{cousinweb}
\textbf{\zh{手心}} (\emph{shǒuxīn}) is built from \zh{手} 'hand' and \zh{心} 'heart', both inherited from Middle Chinese, \emph{syuwX} and \emph{sim} (Guangyun, in Baxter's transcription): the heart, or middle, of the hand.
\end{cousinweb}

\begin{grammarlens}[title={What you've built}]
One more word for saying how it was.
\end{grammarlens}

\subsection*{Your turn}
\begin{itemize}
\raggedright
  \item \emph{Say it:} \emph{shǒuxīn}
  \item \emph{Say it:} \emph{shǒuxīn}, once more
  \item \emph{Say it:} \emph{wǒ shǒuxīn hěn rè shì bu shì tǐwēn gāo le} — My palms are hot. Is my temperature up?
  \item \emph{Recall:} say the Mandarin for to put one's strength into it, then the Mandarin for a finger, then say \emph{shǒuxīn} again
\end{itemize}

\begin{tcolorbox}[breakable,colback=teal!4,colframe=teal!35!black,title={Before you move on}]
What does \zh{手心} mean? (\textquotedblleft{}The palm of the hand\textquotedblright{}.) Say it once more.
\end{tcolorbox}
\section[xīgài]{xīgài — knee}

\label{lesson:ZH-C475-xigai}



\begin{quote}
Before the new one: say the Mandarin for a finger, then the Mandarin for the palm of the hand.
\end{quote}

\subsection*{You'll want to know: xīgài}
\emph{xīgài} — \textquotedblleft{}knee\textquotedblright{}.

\begin{quote}
\emph{wǒ zǒu le yì tiān lù xīgài hěn téng} — I walked all day, and my knees really hurt.
\end{quote}

Its characters come later; for now it is a word for the ear.

\begin{cousinweb}
\emph{xīgài} is built from \emph{xī} 'the knee' and \emph{gài} 'a cover', both inherited from Middle Chinese, \emph{sit} and \emph{kajH} (Guangyun, in Baxter's transcription), which explains the first, after an older book, as the joint of the leg and glosses the second 'to cover'.
\end{cousinweb}

\begin{grammarlens}[title={What you've built}]
One more word for saying how it was.
\end{grammarlens}

\subsection*{Your turn}
\begin{itemize}
\raggedright
  \item \emph{Say it:} \emph{xīgài}
  \item \emph{Say it:} \emph{xīgài}, once more
  \item \emph{Say it:} \emph{wǒ zǒu le yì tiān lù xīgài hěn téng} — I walked all day, and my knees really hurt.
  \item \emph{Recall:} say the Mandarin for a finger, then the Mandarin for the palm of the hand, then say \emph{xīgài} again
\end{itemize}

\begin{tcolorbox}[breakable,colback=teal!4,colframe=teal!35!black,title={Before you move on}]
What does xīgài mean? (\textquotedblleft{}Knee\textquotedblright{}.) Say it once more.
\end{tcolorbox}
\section[hòubèi]{hòubèi — the back of the body}

\label{lesson:ZH-C475-houbei}



\begin{quote}
Before the new one: say the Mandarin for the palm of the hand, then the Mandarin for the knee.
\end{quote}

\subsection*{You'll want to know: hòubèi}
\emph{hòubèi} — \textquotedblleft{}the back of the body\textquotedblright{}.

\begin{quote}
\emph{wǒ gǎnjué hòubèi yǒudiǎnr téng} — My back feels a bit sore.
\end{quote}

Its characters come later; for now it is a word for the ear.

\begin{cousinweb}
\emph{hòubèi} is built from \zh{后} 'behind' and \emph{bèi} 'the back', both inherited from Middle Chinese, \emph{huwX} and \emph{pwojH} (Guangyun, in Baxter's transcription), which glosses the second 'the spine and back'.
\end{cousinweb}

\begin{grammarlens}[title={What you've built}]
One more word for saying how it was.
\end{grammarlens}

\subsection*{Your turn}
\begin{itemize}
\raggedright
  \item \emph{Say it:} \emph{hòubèi}
  \item \emph{Say it:} \emph{hòubèi}, once more
  \item \emph{Say it:} \emph{wǒ gǎnjué hòubèi yǒudiǎnr téng} — My back feels a bit sore.
  \item \emph{Recall:} say the Mandarin for the palm of the hand, then the Mandarin for the knee, then say \emph{hòubèi} again
\end{itemize}

\begin{tcolorbox}[breakable,colback=teal!4,colframe=teal!35!black,title={Before you move on}]
What does hòubèi mean? (\textquotedblleft{}The back of the body\textquotedblright{}.) Say it once more.
\end{tcolorbox}
\section[\texorpdfstring{\zh{身高} (shēngāo)}{身高 (shēngāo)}]{\zh{身高} (shēngāo) — a person's height; body height}

\label{lesson:ZH-C475-shengao}



\begin{quote}
Before the new one: say the Mandarin for the knee, then the Mandarin for the back of the body.
\end{quote}

\subsection*{You'll want to know: \zh{身高}}
\textbf{\zh{身高}} — \emph{shēngāo} — \textquotedblleft{}a person's height; body height\textquotedblright{}.

\begin{quote}
\textbf{\zh{你身高多少}} \emph{nǐ shēngāo duōshao} — How tall are you?
\end{quote}

\begin{cousinweb}
\textbf{\zh{身高}} (\emph{shēngāo}) is built from \zh{身} 'body' and \zh{高} 'high', both inherited from Middle Chinese, \emph{syin} and \emph{kaw} (Guangyun, in Baxter's transcription), which glosses the first 'oneself' and the second 'high'.
\end{cousinweb}

\begin{grammarlens}[title={What you've built}]
One more word for saying how it was.
\end{grammarlens}

\subsection*{Your turn}
\begin{itemize}
\raggedright
  \item \emph{Say it:} \emph{shēngāo}
  \item \emph{Say it:} \emph{shēngāo}, once more
  \item \emph{Say it:} \emph{nǐ shēngāo duōshao} — How tall are you?
  \item \emph{Recall:} say the Mandarin for the knee, then the Mandarin for the back of the body, then say \emph{shēngāo} again
\end{itemize}

\begin{tcolorbox}[breakable,colback=teal!4,colframe=teal!35!black,title={Before you move on}]
What does \zh{身高} mean? (\textquotedblleft{}A person's height; body height\textquotedblright{}.) Say it once more.
\end{tcolorbox}
\section[wān]{wān — to bend}

\label{lesson:ZH-C475-wan-bend}



\begin{quote}
Before the new one: say the Mandarin for the back of the body, then the Mandarin for a person's height.
\end{quote}

\subsection*{You'll want to know: wān}
\emph{wān} — \textquotedblleft{}to bend\textquotedblright{}.

\begin{quote}
\emph{wǒ de gēbo bù néng wān} — I can't bend my arm.
\end{quote}

Its characters come later; for now it is a word for the ear.

\begin{cousinweb}
\emph{wān} is inherited from Middle Chinese \emph{'waen} (Guangyun, in Baxter's transcription), explained after an older book as drawing a bow with the arrow set: the bow bends.
\end{cousinweb}

\begin{grammarlens}[title={What you've built}]
One more word for saying how it was.
\end{grammarlens}

\subsection*{Your turn}
\begin{itemize}
\raggedright
  \item \emph{Say it:} \emph{wān}
  \item \emph{Say it:} \emph{wān}, once more
  \item \emph{Say it:} \emph{wǒ de gēbo bù néng wān} — I can't bend my arm.
  \item \emph{Recall:} say the Mandarin for the back of the body, then the Mandarin for a person's height, then say \emph{wān} again
\end{itemize}

\begin{tcolorbox}[breakable,colback=teal!4,colframe=teal!35!black,title={Before you move on}]
What does wān mean? (\textquotedblleft{}To bend\textquotedblright{}.) Say it once more.
\end{tcolorbox}
\section[(dialogue)]{Looking back: from Faults to Head to Knees}

\label{lesson:ZH-R475-review-470-475}



\begin{quote}
Say the Mandarin for selfish and stubborn.
\end{quote}

\subsection*{Your turn}
\begin{itemize}
\raggedright
  \item \emph{Say it:} selfish, stubborn, full of oneself, slapdash, slow-witted
  \item \emph{Say it:} tears, right before one's eyes, petty and narrow-minded, short-sighted, to wink
  \item \emph{Say it:} an eyebrow, the colour of one's face, to turn hostile suddenly, the look on one's face, the forehead
  \item \emph{Say it:} the mouth, a dentist, the teeth, to have a tooth pulled, to have a tooth filled
  \item \emph{Say it:} the arm, physical strength, stamina, to put one's strength into it, a finger
  \item \emph{Say it:} the palm of the hand, the knee, the back of the body, a person's height, to bend
  \item \emph{Read it:} \zh{眼} \zh{脸} \zh{牙} \zh{力} — each one aloud, then a word that holds it
\end{itemize}

\begin{tcolorbox}[breakable,colback=teal!4,colframe=teal!35!black,title={Before you move on}]
Selfish? (\emph{zìsī}.) Stubborn? (\emph{gùzhí}.) Full of oneself? (\emph{jiāo'ào}.) Slapdash? (\emph{mǎhu}.) Slow-witted? (\emph{bèn}.) Tears? (\emph{yǎnlèi}.) Right before one's eyes? (\textbf{\zh{眼前}}, \emph{yǎnqián}.) Petty and narrow-minded? (\textbf{\zh{小心眼}}, \emph{xiǎoxīnyǎn}.) Short-sighted? (\emph{jìnshì}.) To wink? (\emph{zhǎyǎn}.) An eyebrow? (\emph{méimao}.) The colour of one's face? (\textbf{\zh{脸色}}, \emph{liǎnsè}.) To turn hostile suddenly? (\textbf{\zh{变脸}}, \emph{biànliǎn}.) The look on one's face? (\emph{biǎoqíng}.) The forehead? (\emph{étóu}.) The mouth? (\emph{zuǐ}.) A dentist? (\textbf{\zh{牙医}}, \emph{yáyī}.) The teeth? (\emph{yáchǐ}.) To have a tooth pulled? (\emph{báyá}.) To have a tooth filled? (\emph{bǔyá}.) The arm? (\emph{gēbo}.) Physical strength? (\textbf{\zh{力气}}, \emph{lìqi}.) Stamina? (\textbf{\zh{体力}}, \emph{tǐlì}.) To put one's strength into it? (\textbf{\zh{用力}}, \emph{yònglì}.) A finger? (\emph{shǒuzhǐ}.) The palm of the hand? (\textbf{\zh{手心}}, \emph{shǒuxīn}.) The knee? (\emph{xīgài}.) The back of the body? (\emph{hòubèi}.) A person's height? (\textbf{\zh{身高}}, \emph{shēngāo}.) To bend? (\emph{wān}.)
\end{tcolorbox}
